\documentclass{article}
\usepackage{graphicx} % Required for inserting images
\usepackage{amsfonts, amsmath, amssymb, amsthm, float}
\usepackage{tikz-cd}

\title{Teoria KAM}
\author{Alessandro Ballini}
\date{June 2026}

\begin{document}

\maketitle
\tableofcontents
\newpage
\section{Sistemi hamiltoniani}
Consideriamo una funzione liscia $H$ definita su un aperto $M$ di $\mathbb{T}^n \times \mathbb{R}^n$ = $\{(\theta, r)\}$ (dove $\mathbb{T}^n = \mathbb{R}^n/\mathbb{Z}^n$ è il toro $n$-dimensionale e $(\theta, r) = (\theta_1, \dots, \theta_n, r_1, \dots, r_n)$). \\
Il \textit{campo vettoriale hamiltoniano} $X_H$ indotto da $H$ tramite le equazioni di Hamilton è dato da 
\begin{align} X_H =
    \begin{cases}
        \dot{\theta}_j &= \frac{\partial H}{\partial r_j} \\
        \dot{r}_j &= -\frac{\partial H}{\partial \theta_j}
    \end{cases} \tag{1.1}
\end{align}
Notiamo immediatamente che $\nabla H \cdot X_H = 0$, il gradiente euclideo di $H$ è ortogonale al campo vettoriale hamiltoniano $X_H$. Infatti 
\begin{align}
    \nabla H \cdot X_H &= \sum_{j = 1}^n \left(\frac{\partial H}{\partial \theta_j}\dot{\theta}_j - \frac{\partial H}{\partial r_j}\dot{r}_j\right) \notag \\
    &= \sum_{j = 1}^n \left(\frac{\partial H}{\partial \theta_j}\frac{\partial H}{\partial r_j} - \frac{\partial H}{\partial r_j}\frac{\partial H}{\partial \theta_j}\right) = 0 \tag{1.1}
\end{align}
In particolare, la proiezione sul piano $\{(\theta_j, r_j)\}$ del campo vettoriale hamiltoniano $X_H$ e del gradiente euclideo $\nabla H$ sono ortogonali (ogni termine della somma è nullo). Questo segue dal fatto che $X_H$ è tangente alle curve di livello di $H$ (curve ad energia fissata), mentre $\nabla H$ punta verso la direzione di massima pendenza (positiva) di $H$.

\newpage
\section{Moti quasiperiodici}
Studiamo ora una classe particolare di sistemi hamiltoniani: i sistemi hamiltoniani \textit{integrabili}. Un sistema hamiltoniano integrabile non dipende dalla variabile del moto $\theta$. Pertanto, fissato $r$ si ottiene 
\begin{align}
    \begin{cases}
        \dot{\theta} &= \nabla_rH = cst \\
        \dot{r} &= 0
    \end{cases} \tag{2.1}
\end{align}
dove con $\nabla_rH$ si intende il gradiente rispetto alla coordinata $r = (r_1, \dots, r_n)$, ovvero 
\begin{align}
    \nabla_rH = \left(\frac{\partial H}{\partial r_1}, \dots, \frac{\partial H}{\partial r_n}\right) \notag
\end{align}
Possiamo facilmente determinare il flusso associato a $X_H$ usando la definizione di flusso di un campo vettoriale 
\begin{align}
    \varphi_t(\theta, r) = \left(\theta + t\nabla_rH, r\right) \tag{2.2}
\end{align}
Notiamo che lo spazio delle fasi è costituito da tori invarianti, infatti fissato $r_0 \in \mathbb{R}^n$ si ha 
\begin{align}
    \varphi_t\left(\mathbb{T}^n \times \{r_0\}\right) = \mathbb{T}^n \times \{r_0\} \tag{2.3}
\end{align}
Il flusso ristretto ad uno qualsiasi dei tori invarianti è \textit{quasiperiodico} con periodo determinata dal \textit{vettore frequenza} $\omega = \nabla_rH \in \mathbb{R}^n$. \\ \\
\textbf{Definizione:} sia $M$ una varietà differenziabile e siano $\Phi_1, \Phi_2$ campi vettoriali definiti su $M$. Inoltre, siano $\varphi_1$ e $\varphi_2$ rispettivamente il flusso di $\Phi_1$ e $\Phi_2$. I flussi $\varphi_1$ e $\varphi_2$ vengono detti \textit{topologicamente coniugati} se esiste un omeomorfismo $h$ su $M$ tale che
\begin{align}
    h(\varphi_1(t, x)) = \varphi_2(t, h(x)) \tag{2.4}
\end{align}
Notiamo che la coniugazione topologica è una condizione più forte dell'equivalenza topologica, infatti si richiede che venga rispettato anche l'ordine temporale. \\ \\
\textbf{Definizione:} un'applicazione continua $p : R \longrightarrow X$ di spazi topologici è detta rivestimento se è suriettiva e se per ogni $x \in X$ ammette un intorno aperto $U_x$ tale che $p^{-1}(U_x)$ è unione disgiunta di aperti di $R$ e $p^{-1}(U_x)$ ognuno dei quali viene mappato omeomorficamente su $U_x$. \\ \\
\textbf{Definizione:} sia $f : X \longrightarrow Y$ un'applicazione continua di spazi topologici 
\begin{align}
    p : \tilde{X} \longrightarrow X \notag \\
    q : \tilde{Y} \longrightarrow Y \notag
\end{align}
rivestimenti. $\tilde{f} : \tilde{X} \longrightarrow \tilde{Y}$ è detto \textit{sollevamento} di $f$ se $q \circ \tilde{f} = f \circ p$. \\ \\
\textbf{Osservazione:} Notiamo che il sollevamento di un'applicazione continua non è unico, infatti, in base alle proprietà del rivestimento si possono trovare più sollevamenti. La non unicità dipende dal fatto che ine generale $p$ non è iniettiva, quindi $p^{-1}(x)$ può contenere più di un punto. \\ \\
Dato $X$ spazio topologico e $x_0 \in X$ chiamiamo spazio topologico \textit{puntato} o \textit{con punto base} la coppia $(X, x_0)$. Inoltre, dati due spazi topologici puntati $(X, x_0), (Y, y_0)$ chiamiamo \textit{applicazione continua di spazi topologici} un'applicazione continua $f : (X, x_0) \longrightarrow (Y, y_0)$ tale che $f(x_0) = y_0$. \\ \\
\textbf{Teorema (di unicità del sollevamento):} siano $p : (R, r_0) \longrightarrow (Y, y_0)$ un rivestiemnto e $f : (X, x_0) \longrightarrow (Y, y_0)$ un'applicazione continua di spazi topologici puntati con $Y$ connesso. Se esiste $\tilde{f} : (X, x_0) \longrightarrow (R, r_0)$ continua tale che $p \circ \tilde{f} = f$, è unica. \\

\textbf{Dimostrazione:} supponiamo che esista $\tilde{f} : (X, x_0) \longrightarrow (R, r_0)$ tale che $p \circ \tilde{f} = f$. Per definizione di applicazione di spazi topologici puntati si ha $p(r_0) = y_0$, $f(x_0) = y_0$ e $\tilde{f}(x_0) = r_0$. Inoltre, per definizione di rivestimento, esiste un intorno aperto $U_{y_0}$ di $y_0$ tale che $V_{y_0} := p^{-1}(U_{y_0})$ è unione disgiunta di aperti di $R$ e viene mappato omeomorficamente su $U_{y_0}$ da $p$.   \\ \\
Fissiamo $r \in \mathbb{R}^n$ e consideriamo il flusso 
\begin{align}
    \varphi_t : &\mathbb{T}^n \longrightarrow \mathbb{T}^n \qquad \theta \mapsto \theta + t\alpha \notag
\end{align}
dove $\alpha = \nabla_rH(r)$. \\ \\
\textbf{Lemma 2.1:} il vettore frequenza $\alpha$ è invariante per coniugazione topologica a meno di azione di elementi di $GL_n(\mathbb{Z})$: se due flussi $\theta \mapsto \theta + t\alpha$ e $\theta \mapsto t\beta$ con $\alpha, \beta \in \mathbb{R}^n$ sono topologicamente coniugati, esiste $A \in GL_n(\mathbb{Z})$ tale che $\alpha = A\beta$. In particolare se l'omeomorfismo che li allaccia preserva l'orientamento si ha $A \in SL_n(\mathbb{Z})$. \\

\textbf{Dimostrazione:} supponiamo che i flussi definiti da $\theta + t\alpha$ e $\theta + t\beta$ siano topologicamente coniugati. Per definizione esiste un omeomorfismo $h$ di $\mathbb{T}^n$ tale che 
\begin{align}
    h(\theta + t\alpha) = h(\theta) + t\beta \tag{2.5}
\end{align}
A meno di sostituire $h(\theta)$ con $h(\theta) - h(0)$ possiamo supporre $h(0) = 0$. Sia $H : \mathbb{R}^n \longrightarrow \mathbb{R}^n$ l'unico sollevamento di $h$ tale che $H(0) = 0$.


\end{document}